\documentclass[10pt,a4paper]{article}
\usepackage[utf8]{inputenc}
\usepackage[T1]{fontenc}
\usepackage[margin=20mm]{geometry}
\usepackage{array,booktabs,tabularx}
\usepackage{amsmath,amssymb}
\usepackage{xcolor}
\usepackage[colorlinks=true,linkcolor=blue!60!black,citecolor=green!45!black,urlcolor=blue!65!black]{hyperref}

\pagestyle{empty}
\setlength{\parindent}{0pt}
\setlength{\parskip}{0.22em}
\setlength{\tabcolsep}{3.2pt}
\renewcommand{\arraystretch}{1.08}
\makeatletter
\renewcommand\section{\@startsection{section}{1}{0pt}{0.65ex plus .2ex}{0.35ex}{\normalfont\normalsize\bfseries}}
\makeatother

\begin{document}
\small

\begin{center}
{\large\bfseries Condensed Novelty Check: Delay--Layout Coupling in UWB Self-Calibration}\\[-0.1em]
{\small Executive summary for paper planning}
\end{center}

\section{Thesis}

UWB anchor self-calibration can trade antenna-delay parameters against layout scale, so same-environment tag accuracy is not a reliable proxy for physical calibration quality. In the Erlangen campaign, V4 is metrically distorted (Sim3 scale 0.958) and V5 is metric-correct (1.010), yet V4 wins under p50 static LOO positioning (57.9 vs 67.8~mm); after reducing static tag-anchor positive range bias, V5 wins (44.5 vs 53.9~mm). This ranking flip is the key evidence that the wrong geometry can win by cancelling structured NLOS bias.

\section{Literature Gap}

\begin{center}
\scriptsize
\begin{tabularx}{\textwidth}{@{}>{\raggedright\arraybackslash}X>{\raggedright\arraybackslash}X@{}}
\toprule
What is known & What is NOT known (our gap) \\
\midrule
UWB self-calibration estimates anchors from radio ranges, robot motion, multihop links, or auxiliary sensing \cite{hamer_2018_self_calibrating_ultra_wideband_network,corbalan_2023_self_localization_uwb_anchors,herbruggen_2023_multihop_self_calibration_algori}. & These papers do not test whether the layout with the best tag accuracy is metrically correct. \\
Antenna-delay and range-bias calibration are established UWB problems \cite{ledergerber_2018_calibrating_away_inaccuracies_uwb,de_preter_2019_range_bias_modeling_autocalibration,shalaby_2023_calibration_and_uncertainty_char,liu_2024_data_driven_antenna_delay_calibr}. & Delay is treated as a nuisance correction, not as a parameter that can hide global layout-scale error. \\
NLOS mitigation uses channel statistics, learned models, robust weighting, or biased-range models \cite{wymeersch_2012_a_machine_learning_approach,difranco_2017_calibration_free_network_localization_nlos_uwb}. & NLOS reduction is not used as a causal intervention to reveal which calibration geometry is physically valid. \\
Fisher/CRLB and rigidity analyses describe weak directions and rank deficiencies in localization networks \cite{shen_2010_fundamental_limits_cooperative,aspnes_2006_theory_network_localization}. & Prior work does not connect a weak delay--layout direction to a measured wrong-geometry-wins result in UWB self-calibration. \\
Scale--propagation-parameter ambiguity is known in acoustic self-calibration and related ranging systems \cite{burgess_2015_toa_self_calibration}. & The interaction with structured UWB NLOS bias, and the resulting ranking flip between distorted and metric-correct layouts, is not reported. \\
AniTrack observed that self-localized anchors outperformed surveyed anchors in tag accuracy \cite{luder_2025_anitrack}. & It reports the paradox but does not diagnose Sim3 scale, delay--layout coupling, NLOS cancellation, or profile/Fisher structure. \\
\bottomrule
\end{tabularx}
\end{center}

\section{Our Contribution}

The first contribution is scientific: we quantitatively characterize delay--layout--NLOS coupling in UWB anchor self-calibration. The evidence chain combines the V4/V5 Sim3 split, Fisher information with weakest eigenvalue $10^{-6}$, profile-likelihood and morph-valley surfaces, phase-center robustness, and the raw-frame lower-tail intervention. The decisive result is the ranking flip: V4 wins when p50 tag-anchor positive bias is present, but V5 wins after that bias is reduced, showing that V4's p50 advantage came from cancellation rather than better physical calibration.

The second contribution is methodological: we provide a falsification protocol that separates metric correctness from same-environment positioning accuracy. It uses Sim3 anchor-scale validation, a layout/delay transfer matrix, winner's-curse correction, nested spatial CV, phase-center sensitivity, NLOS leakage auditing, and ROTO dynamic checks to prevent a single accuracy number from hiding the mechanism.

The same failure mode can occur beyond UWB whenever per-device bias and geometric scale share a near-degenerate Fisher direction, including GNSS clock--geometry coupling, acoustic sound-speed--geometry coupling, and other range-only calibration systems.

\section{Nearest Prior Art}

\begin{center}
\scriptsize
\begin{tabularx}{\textwidth}{@{}>{\raggedright\arraybackslash}p{28mm}>{\raggedright\arraybackslash}X>{\raggedright\arraybackslash}X@{}}
\toprule
Paper & What they did & How we differ \\
\midrule
Hamer and D'Andrea 2018 \cite{hamer_2018_self_calibrating_ultra_wideband_network} & Built a self-calibrating one-way UWB network with clock synchronization and anchor computation for multi-robot localization. & We test whether the empirically best radio-only layout is physically correct, and show that a distorted layout can win by cancelling NLOS bias. \\
Lutz et al. 2019 \cite{lutz_2019_visual_inertial_slam_uwb_error_model} & Used visual-inertial SLAM and fiducial markers to estimate UWB anchor poses and sensor error-model parameters. & Their calibration is aided by external visual structure; ours diagnoses a range-only self-calibration degeneracy using Vicon only as independent validation. \\
Ledergerber and D'Andrea 2018 \cite{ledergerber_2018_calibrating_away_inaccuracies_uwb} & Estimated module-dependent range inaccuracies by maximum likelihood without external ground truth. & We show that fitting delay/range inaccuracies can reward a geometrically wrong solution unless anchor metric scale is checked separately. \\
Luder et al. 2025 (AniTrack) \cite{luder_2025_anitrack} & Reported that self-localized DWM3000 anchors gave better tag accuracy than laser-surveyed anchors (13.96 vs 16.57~cm average error). & They report an unexplained paradox; we provide the mechanism, Sim3/Fisher/profile evidence, and the NLOS-reduction intervention. \\
Pan et al. 2015 (GNSS) \cite{pan_2015_realtime_ppp_clock_orbit} & Showed satellite-clock estimation can absorb orbit error and improve PPP positioning accuracy. & Their coupled-error absorption is a GNSS clock/orbit effect; our result is a UWB self-calibrated-geometry error that makes the physically wrong layout beat the correct one. \\
\bottomrule
\end{tabularx}
\end{center}

\section{Risk and Status}

\begin{center}
\scriptsize
\begin{tabularx}{\textwidth}{@{}>{\raggedright\arraybackslash}p{34mm}>{\raggedright\arraybackslash}p{19mm}>{\raggedright\arraybackslash}X@{}}
\toprule
Risk & Severity & One-line mitigation \\
\midrule
One room and 24 static positions & High & State that transfer remains unproven and freeze the pipeline for second-room validation. \\
Cross-domain scale--delay ambiguity is prior art & Medium & Claim quantitative UWB characterization plus NLOS-induced ranking flip, not the bare ambiguity. \\
Lower-tail result is static batch, not real-time & Medium & Present it as a diagnostic intervention and static engineering candidate, not a deployable NLOS classifier. \\
High-error tail remains & Medium & Separate median mechanism evidence from P95/dynamic performance claims. \\
Phase-center uncertainty & Low & Use sensitivity analysis and avoid treating Vicon-oracle rank as the sole proof. \\
\bottomrule
\end{tabularx}
\end{center}

Status: 129 papers indexed, 15 cited, 0 explain delay--layout--NLOS coupling. Target: \emph{Measurement} (IF $\sim$5.6). Title: ``When Wrong Geometry Wins: Delay--Layout Coupling in UWB Self-Calibration.''

\begingroup
\scriptsize
\begin{thebibliography}{99}
\bibitem{aspnes_2006_theory_network_localization} J. Aspnes, T. Eren, D. K. Goldenberg, A. S. Morse, W. Whiteley, Y. R. Yang, B. D. O. Anderson, and P. N. Belhumeur. A Theory of Network Localization. \emph{IEEE Transactions on Mobile Computing}, 5(12):1663--1678, 2006.
\bibitem{burgess_2015_toa_self_calibration} S. Burgess, Y. Kuang, and K. \AA{}str\"om. TOA Sensor Network Self-Calibration for Receiver and Transmitter Spaces with Difference in Dimension. \emph{Signal Processing}, 107:33--42, 2015.
\bibitem{corbalan_2023_self_localization_uwb_anchors} P. Corbalan, G. P. Picco, M. Coors, and V. Jain. Self-Localization of Ultra-Wideband Anchors: From Theory to Practice. \emph{IEEE Access}, 11:29711--29725, 2023.
\bibitem{difranco_2017_calibration_free_network_localization_nlos_uwb} C. Di Franco, A. Prorok, N. Atanasov, B. Kempke, P. Dutta, V. Kumar, and G. J. Pappas. Calibration-Free Network Localization using Non-Line-of-Sight UWB Measurements. \emph{IPSN}, 2017.
\bibitem{hamer_2018_self_calibrating_ultra_wideband_network} M. Hamer and R. D'Andrea. Self-Calibrating Ultra-Wideband Network Supporting Multi-Robot Localization. \emph{IEEE Access}, 6:22292--22304, 2018.
\bibitem{herbruggen_2023_multihop_self_calibration_algori} B. Van Herbruggen, S. Luchie, J. Fontaine, and E. De Poorter. Multihop Self-Calibration Algorithm for Ultra-Wideband Anchor Node Positioning. \emph{IEEE Journal of Indoor and Seamless Positioning and Navigation}, 2023.
\bibitem{ledergerber_2018_calibrating_away_inaccuracies_uwb} A. Ledergerber and R. D'Andrea. Calibrating Away Inaccuracies in Ultra Wideband Range Measurements: A Maximum Likelihood Approach. \emph{IEEE Access}, 6:78719--78730, 2018.
\bibitem{liu_2024_data_driven_antenna_delay_calibr} Z. Liu, T. Hakala, J. Hyyppa, A. Kukko, H. Kaartinen, and R. Chen. Data-Driven Antenna Delay Calibration for UWB Devices for Network Positioning. \emph{IEEE Transactions on Instrumentation and Measurement}, 2024.
\bibitem{lutz_2019_visual_inertial_slam_uwb_error_model} P. Lutz, M. J. Schuster, and F. Steidle. Visual-inertial SLAM aided estimation of anchor poses and sensor error model parameters of UWB radio modules. \emph{ICAR}, 2019.
\bibitem{luder_2025_anitrack} V. Luder, L. Schulthess, S. Cortesi, L. Rivero Davis, and M. Magno. AniTrack: A Power-Efficient, Time-Slotted and Robust UWB Localization System for Animal Tracking in a Controlled Setting. \emph{arXiv:2506.00216}, 2025.
\bibitem{pan_2015_realtime_ppp_clock_orbit} S. Pan, W. Chen, X. Jin, X. Shi, and F. He. Real-Time PPP Based on the Coupling Estimation of Clock Bias and Orbit Error with Broadcast Ephemeris. \emph{Sensors}, 15(7):17808--17826, 2015.
\bibitem{shalaby_2023_calibration_and_uncertainty_char} M. Shalaby, C. C. Cossette, J. R. Forbes, and J. Le Ny. Calibration and Uncertainty Characterization for Ultra-Wideband Two-Way-Ranging Measurements. \emph{ICRA}, 2023.
\bibitem{shen_2010_fundamental_limits_cooperative} Y. Shen, H. Wymeersch, and M. Z. Win. Fundamental Limits of Wideband Localization--Part II: Cooperative Networks. \emph{IEEE Transactions on Information Theory}, 56(10):4981--5000, 2010.
\bibitem{wymeersch_2012_a_machine_learning_approach} H. Wymeersch, S. Marano, W. M. Gifford, and M. Z. Win. A Machine Learning Approach to Ranging Error Mitigation for UWB Localization. \emph{IEEE Transactions on Communications}, 2012.
\bibitem{de_preter_2019_range_bias_modeling_autocalibration} A. De Preter, G. Goysens, J. Anthonis, J. Swevers, and G. Pipeleers. Range Bias Modeling and Autocalibration of an UWB Positioning System. \emph{IPIN}, 2019.
\end{thebibliography}
\endgroup

\end{document}
